%% ---------------------------------------------------------------------------
%% Identifiability of OMCOA and MCOA in the noiseless case
%% (uses the macros of the manuscript: \X, \U, \V, \uu, \vv, \Stiefel, \diag)
%% ---------------------------------------------------------------------------

\begin{proposition}[Noiseless identifiability]\label{prop:identif}
Let $\X_k=\U\Sigma_k\V_k^\intercal$, $k=1,\ldots,K$, with $\U\in\Stiefel{N}{R}$,
$\V_k\in\Stiefel{p_k}{R}$, $\Sigma_k=\diag(\sigma_k^{(1)},\ldots,\sigma_k^{(R)})$,
$\sigma_k^{(r)}\geq 0$, and assume that every dimension is present in at least
one block: for each $r$ there is a $k$ with $\sigma_k^{(r)}>0$. Denote by
$\mathbf s^{(r)}=(\sigma_1^{(r)},\ldots,\sigma_K^{(r)})$ the profile of dimension
$r$ across the blocks and, for block weights $w_1,\ldots,w_K>0$, by
\[
  \lambda_r(\mathbf w)=\sum_{k=1}^K w_k\bigl(\sigma_k^{(r)}\bigr)^2
\]
the squared singular values of the weighted concatenated matrix
$[\sqrt{w_1}\X_1|\cdots|\sqrt{w_K}\X_K]$.
\begin{enumerate}
  \item[(a)] \textbf{OMCOA.} If the profiles are pairwise distinct,
  $\mathbf s^{(r)}\neq\mathbf s^{(r')}$ for $r\neq r'$, then every minimiser of
  criterion~\eqref{Crit} in dimension $R$, weighted or not, is equal to
  $(\U,\V_k,\Sigma_k)$ up to a common permutation of the dimensions and to
  the joint change of sign $(\uu^{(r)},\vv_1^{(r)},\ldots,\vv_K^{(r)})\mapsto
  -(\uu^{(r)},\vv_1^{(r)},\ldots,\vv_K^{(r)})$; the columns $\vv_k^{(r)}$ with
  $\sigma_k^{(r)}=0$ are not determined.
  \item[(b)] \textbf{MCOA.} If the values $\lambda_1(\mathbf w),\ldots,
  \lambda_R(\mathbf w)$ are pairwise distinct, MCOA applied to the weighted
  blocks $\sqrt{w_k}\X_k$ recovers $(\U,\V_k)$ up to the order and the signs of
  the dimensions, and the weights $\sqrt{w_k}\,\sigma_k^{(r)}$. Its solution
  then coincides with the OMCOA solution (once the MCOA weights are divided by
  $\sqrt{w_k}$).
  \item[(c)] If $\lambda_r(\mathbf w)=\lambda_{r'}(\mathbf w)$ for some
  $r\neq r'$, the MCOA solution is not determined by the data: at the first
  step where the largest remaining value is tied, any unit vector of the
  corresponding eigenspace is a solution, and the computed solution differs
  in general from $(\U,\V_k)$, even when the profiles are distinct.
\end{enumerate}
\end{proposition}

Since equal profiles give equal values $\lambda_r(\mathbf w)$ for every
$\mathbf w$, the condition of (b) implies that of (a), whatever the weights.
Conversely, for any weights there are designs with distinct profiles and tied
values $\lambda_r(\mathbf w)$: OMCOA identifies strictly more configurations
than MCOA, and its identifiability does not depend on the block weighting.

\begin{proof}
\emph{(a)} The true parameters give the value $0$ to criterion~\eqref{Crit}, so
every minimiser $(\U',\V_k',\Sigma_k')$ satisfies
$\X_k=\U'\Sigma_k'\V_k'^\intercal$ for all $k$; this does not depend on the
weights given to the blocks. Since every dimension is present in some block,
the column space of $\X=[\X_1|\cdots|\X_K]$ is the span of $\U$, of dimension
$R$. It is contained in the span of $\U'$, which also has dimension $R$; hence
$\U'=\U\mathbf O$ with $\mathbf O$ orthogonal. From
$\X_k\X_k^\intercal=\U\Sigma_k^2\U^\intercal=\U'\Sigma_k'^2\U'^\intercal$ we get
\[
  \mathbf O^\intercal\Sigma_k^2\,\mathbf O=\Sigma_k'^2,\qquad k=1,\ldots,K .
\]
For $\mathbf c\in\mathbb R^K$, the diagonal matrix
$\mathbf A(\mathbf c)=\sum_k c_k\Sigma_k^2$ has entries
$a_r=\sum_k c_k(\sigma_k^{(r)})^2$, and $\mathbf O^\intercal\mathbf A(\mathbf c)\mathbf O$
is diagonal. Since the profiles are distinct and nonnegative, the squared
profiles are distinct, and $\mathbf c$ can be chosen outside the finite union
of hyperplanes $\{a_r=a_{r'}\}$, so that the entries of $\mathbf A(\mathbf c)$
are pairwise distinct. An orthogonal matrix that diagonalises a diagonal
matrix with distinct entries is a signed permutation matrix:
$\mathbf O=\mathbf P\mathbf D$, with $\mathbf P$ a permutation matrix and
$\mathbf D=\diag(\pm1)$. Hence $\U'=\U\mathbf P\mathbf D$ and
$\Sigma_k'^2=\mathbf P^\intercal\Sigma_k^2\mathbf P$, that is
$\Sigma_k'=\mathbf P^\intercal\Sigma_k\mathbf P$ since the weights are
nonnegative. Finally, $\U'^\intercal\X_k=\Sigma_k'\V_k'^\intercal$ gives, for
every $(k,r)$ with $\sigma_k'^{(r)}>0$,
$\vv_k'^{(r)}=\X_k^\intercal\uu'^{(r)}/\sigma_k'^{(r)}$, which is the matching
column of $\V_k$ with the sign of $\uu'^{(r)}$; the columns with
$\sigma_k'^{(r)}=0$ are only constrained by $\V_k'\in\Stiefel{p_k}{R}$.

\emph{(b)} Write $\mathcal R_t$ for the set of dimensions not yet extracted at
step $t$ of MCOA ($\mathcal R_1=\{1,\ldots,R\}$), and assume that the deflated
blocks are
$\X_k^{(t)}=\sqrt{w_k}\sum_{r\in\mathcal R_t}\sigma_k^{(r)}\uu^{(r)}\vv_k^{(r)\intercal}$,
which holds for $t=1$. Then
\[
  \X^{(t)}\X^{(t)\intercal}=\sum_{k}\X_k^{(t)}\X_k^{(t)\intercal}
  =\sum_{r\in\mathcal R_t}\lambda_r(\mathbf w)\,\uu^{(r)}\uu^{(r)\intercal},
\]
because the columns of each $\V_k$ are orthonormal. The values
$\lambda_r(\mathbf w)$ being distinct, the leading left singular vector of
$\X^{(t)}$ is $\pm\uu^{(r_t)}$, with $r_t=\argmax_{r\in\mathcal R_t}\lambda_r(\mathbf w)$.
Then $\X_k^{(t)\intercal}\uu^{(r_t)}=\sqrt{w_k}\,\sigma_k^{(r_t)}\vv_k^{(r_t)}$:
the MCOA weight is $\sqrt{w_k}\,\sigma_k^{(r_t)}$ and, when it is positive, the
normalised loading is $\pm\vv_k^{(r_t)}$ (with the sign of $\uu$), so that the
deflation $\X_k^{(t)}(\mathbf I-\vv_k^{(r_t)}\vv_k^{(r_t)\intercal})$ removes exactly the
term $r_t$. When the weight is zero, the term $r_t$ is absent from block $k$
and the block is left unchanged. The induction hypothesis holds at step
$t+1$ with $\mathcal R_{t+1}=\mathcal R_t\setminus\{r_t\}$. After $R$ steps,
MCOA has recovered $\U$ and the $\V_k$ (columns with a positive weight) up to
order and signs, and gives a zero value to criterion~\eqref{Crit} for the
weighted blocks: by (a), it coincides with the OMCOA solution.

\emph{(c)} Let $t$ be the first step where
$\max_{r\in\mathcal R_t}\lambda_r(\mathbf w)$ is attained by several dimensions
$r\in\mathcal T\subset\mathcal R_t$. Up to step $t-1$, the argument of (b)
applies, and the leading eigenvalue of $\X^{(t)}\X^{(t)\intercal}$ has the
eigenspace $\mathrm{span}\{\uu^{(r)},r\in\mathcal T\}$, of dimension
$|\mathcal T|\geq2$. Every unit vector of this space maximises the MCOA
criterion at step $t$; a vector other than $\pm\uu^{(r)}$, $r\in\mathcal T$,
gives a common component that is not a column of $\U$ up to sign, hence, by
(a), a solution different from the OMCOA one.
\end{proof}

\begin{remark}
In \texttt{ade4}, the weight $w_k$ combines the column weights of the
\texttt{ktab} object (by default $1/p_k$ for each column of block $k$ with
\texttt{ktab.list.df}) and the block weighting of \texttt{mcoa}: $1$ for
\texttt{option="uniform"}, the inverse of the total inertia for
\texttt{"inertia"}, the inverse of the first eigenvalue for \texttt{"lambda1"}.
The condition of (b) therefore depends on these choices: for instance, with
blocks of different sizes, the default column weights $1/p_k$ can break a tie
of the unweighted values $\sum_k(\sigma_k^{(r)})^2$, and conversely.
Standardising the columns of the blocks (\texttt{dudi.pca(scale=TRUE)}) is not
a block weighting: it multiplies $\X_k$ on the right by a diagonal matrix
$\mathbf D_k^{-1}$, and $\X_k\mathbf D_k^{-1}$ is no longer of the form
$\U\Sigma_k\V_k^\intercal$ with orthonormal $\V_k$, unless all the columns of
block $k$ have the same variance.
\end{remark}

\begin{remark}
With noise, exact ties become near ties. When two values $\lambda_r(\mathbf w)$
are close, the corresponding MCOA axes are only determined up to a rotation
within their plane: by the Davis--Kahan theorem \cite{davis1970}, the angle
between estimated and true axes is bounded by the noise level divided by the
gap between the two values, and this bound is attained in the worst case.
OMCOA does not rely on this gap, but on the separation of the profiles. The
simulation with the parameter $t$, which moves continuously from distinct to
tied values of $\lambda_r(\mathbf w)$ with distinct profiles, illustrates the
difference.
\end{remark}

%% BibTeX entry to add to Biblio.bib:
%% @article{davis1970,
%%   author  = {Chandler Davis and W. M. Kahan},
%%   title   = {The Rotation of Eigenvectors by a Perturbation. {III}},
%%   journal = {SIAM Journal on Numerical Analysis},
%%   volume  = {7},
%%   number  = {1},
%%   pages   = {1--46},
%%   year    = {1970},
%%   doi     = {10.1137/0707001}
%% }
